\chapter{Results and Discussion}\label{chap:results}

This chapter presents the findings from applying transcript-level COTAN to scRNA-seq data. It begins with quality assessment (GDI clustering validation) and a comparison of gene-level and transcript-level structure, before presenting the main contribution: DTU detection results across the human and mouse datasets. Each section builds toward answering whether sparse 3'-biased data can support isoform-switch discovery using contingency-table statistics.

\section{Dataset Quality Assessment via GDI}

Before DTU calls can be interpreted, the clustering must be shown to capture biological structure rather than technical noise. The Global Differentiation Index (GDI), defined in Section~\ref{sec:cotan-framework}, provides this validation.

\subsection*{Transcript-level Clustering Quality}\label{sec:gdi-transcript}

Figure~\ref{fig:gditranscript} shows the GDI plot for the mouse dataset clustered on transcript counts. Each point represents a cell, colored by its assigned cluster (11 uniform clusters total). The clear spatial separation indicates that transcript-level co-expression patterns distinguish distinct cellular populations.

Key observations from this visualization:
\begin{itemize}
    \item Cluster boundaries are well-defined with minimal overlap, suggesting robust isoform-based signatures.
    \item High-GDI genes (marker candidates) show selective enrichment in specific clusters rather than uniform expression.
    \item The noisy remainder cells (not assigned to any cluster) form a tight central cloud; these were filtered out from downstream DTU analysis.
\end{itemize}


This confirms that transcript-level data contains sufficient signal for meaningful clustering, despite sparsity and 3'-bias constraints.

\begin{figure}[htbp]
    \centering
    \includegraphics[width=0.8\textwidth]{GDI_transcript_plot_markers.png}
    \caption{GDI plot of the mouse dataset clustered on transcript-level counts, coloured by the 11 uniform clusters.}
    \label{fig:gditranscript}
\end{figure}

\subsection*{Gene vs Transcript Structure Comparison}

Figure~\ref{fig:gdi_genes} shows the same cells clustered using standard gene counts (15 clusters). The gene run reads the independent GEO gene matrix---28,692 features against the transcript matrix's 12,848---so its larger cluster count reflects the denser matrix, not a finer gene-level resolution.

\begin{figure}[htbp]
    \centering
    \includegraphics[width=0.8\textwidth]{GDI_genes_plot_markers.png}
    \caption{GDI plot of the same cells clustered on gene-level counts, yielding 15 clusters.}
    \label{fig:gdi_genes}
\end{figure}

These GDI analyses establish that both clustering approaches capture real biological signal, but with different sensitivity to cellular substructure. The main question is whether this structural difference translates into distinct DTU detection.

\section{Clustering Comparison: Gene vs Transcript Level}\label{sec:clustering-comparison}

To quantify how gene and transcript clustering diverge, a confusion matrix was constructed across 4,077 common cells present in both analyses.

\begin{figure}[htbp]
    \centering
    \includegraphics[width=0.8\textwidth]{heatmap_genes_vs_transcripts.pdf}
    \caption{Confusion matrix comparing transcript-level (11 clusters) and gene-level (15 clusters) assignments for the 4,077 common cells of the mouse dataset.}
    \label{fig:confusion-matrix}
\end{figure}

Figure~\ref{fig:confusion-matrix} shows the confusion matrix heatmap. Several patterns emerge:
\begin{itemize}
    \item \textbf{Diagonal dominance}: Most cells receive consistent assignments across both methods, validating that neither approach produces random clusters.
    
    \item \textbf{Systematic splitting}: Several transcript-level clusters correspond to sub-regions of single gene-level clusters. For example, transcript cluster 4 splits from the broader gene cluster 7, suggesting isoform variation within what appears homogeneous at gene resolution.
    
    \item \textbf{Selective merging}: Conversely, some adjacent transcript clusters merge in gene analysis, indicating that their differences rely on isoform-specific expression patterns rather than overall transcriptional activity.
\end{itemize}

The interpretation: across the two runs the counts differ (11 against 15) because the partitions come from different matrices---the transcript matrix is a sparser, feature-filtered view, the gene matrix a denser GEO count set---so the totals are not a resolution ranking. What the confusion matrix actually shows is agreement in the large and divergence in the particular: the dominant diagonal confirms both partitions recover the same broad populations, while the local splits and merges mark where isoform information changes the grouping. That is precisely the regime where DTU detection should add value---if isoform switches distinguish cell states that genes alone cannot separate, then transcript-aware methods capture biologically relevant signals.

However, this raises a critical question: whether the detected DTU events hold across clustering strategies, or depend entirely on how cells are partitioned. The robustness analysis below answers it.

\section{Core Finding: DTU Detection Results}\label{sec:dtu-results}
The algorithm from Section~\ref{sec:methods-dtu-algorithm} was applied to both human and mouse datasets, filtering for mutually exclusive isoform pairs with bidirectional cluster preference (threshold $\tau$ as specified). The results show that contingency-table statistics can recover within-gene exclusivity patterns even in sparse 3'-biased data.

\subsection*{Human Dataset Findings (Arigoni)}\label{sec:results-human}

The human cell line mixture provided ground-truth validation through known biology. Of the 4958 genes carrying at least two detected transcripts, 430 retain at least one significant mutually exclusive pair\footnote{Counted by \texttt{extract\_dtu\_candidates()} and recorded in the run log (\texttt{results/arrigoni/logs/}); the exported \texttt{dtu\_candidates.csv} holds only the 21 survivors.}, yet after the bidirectional contrast threshold $\tau = 0.2$ only 21 candidates survive, across just 9 genes. That collapse from 430 to 9 is the algorithm working as designed: mutual exclusivity alone admits hundreds of pairs, but only a handful also flip their enrichment between clusters, and it is this second condition---a pair must be preferred in one cell type and disfavoured in another---that separates a genuine switch from a statistical curiosity.

The surviving candidates cohere into a single story. TPM1, the strongest call, is a tropomyosin whose isoforms oppose each other across the cell lines: one isoform is enriched in PC9 and DV90 while the other dominates A549 and HCC78, a reciprocal preference that reads as cell-type-specific cytoskeletal tuning. Alongside it sit the ribosomal proteins RPL10 and RPS24, whose unexpected DTU signal points toward cell-type-specific translation regulation---consistent with the heterogeneity of ribosomal composition underlying specialised translation \cite{xue2012specialized}. The two themes are biologically distinct but converge in method: in both cases a single gene carries two isoforms that the same cell types use in opposite directions, and it is that convergence, not any single gene, that validates the detector. The final list is small precisely because the algorithm demands both \gls{mutualExclusivity} and a bidirectional contrast, and that conservative design is what lets nine genes carry the whole human result.

\subsection*{Mouse Dataset Findings (Ding)}\label{sec:results-mouse}

For mouse brain tissue a lower threshold, $\tau = 0.05$, was applied, and the balance shifts: only 14 genes carry significant pairs, yet 46 candidates survive across 8 genes. The inversion is itself informative---fewer initial pairs but more survivors means mouse cortex switches are gradual rather than sharp, so a relaxed threshold admits many weak but real transitions that the strict human cutoff would have discarded.

The candidates trace an arc across the tissue. Meg3, a long non-coding RNA\defn{Long non-coding RNA (lncRNA)}{A long RNA molecule that does not encode a protein but regulates gene expression, splicing, and development.} with a documented role in neuronal differentiation, switches along the differentiation trajectory itself: isoform A dominates early progenitor clusters while isoform B is enriched in differentiated neurons---a transition that tracks a continuous gradient rather than a binary state. Malat1, another highly conserved lncRNA involved in splicing regulation, shows the same cluster-specific preference, consistent with the alternative 3$'$ isoform usage documented for this locus \cite{ghanam2023malat1}. The two lncRNAs are neighbours in function---both non-coding, both switching with differentiation---and together they anchor the mouse result in established biology. The arc then turns inward: Srrm2 encodes a core spliceosome\defn{Spliceosome}{The molecular machine that removes introns and joins exons during splicing.} component, and its own DTU suggests feedback regulation, alternative splicing of splicing factors being a known mechanism for tuning isoform landscapes across cell states. Finally, Ttc14, a novel candidate with no prior DTU literature, requires \gls{orthogonalValidation} to confirm its biological significance---and its presence is a reminder that the method also surfaces calls no published study has anticipated.

\subsection*{Robustness Analysis: Shared vs Exclusive Events}\label{sec:robustness}

DTU detection was re-run using gene-derived cluster labels on the same transcript matrix, and the resulting event lists were compared across both clustering strategies. Both partitionings returned the same total, 46 DTU pairs; their composition differed:
\begin{itemize}
    \item \textbf{Shared events}: Majority of pairs detected by both methods, including Meg3 and Malat1---these represent strong signals surviving regardless of partitioning scheme.
    
    \item \textbf{Transcript-exclusive (3 pairs)}: Srrm2 and Ttc14 isoform switches appeared only with transcript-derived clusters. These likely reflect subtle transitions requiring fine-grained population resolution to detect.

    \item \textbf{Gene-exclusive (4 pairs)}: Snrnp70, Lrrc7, Zfp941, and a specific Meg3 pair detected only using gene-level clustering. These may represent weaker signals that benefit from aggregated statistical power.
\end{itemize}

The interpretation:
\begin{enumerate}
    \item Strongest DTU signals (high effect size) survive both approaches---validating the algorithm's core functionality.
    
    \item Subtle switches depend on clustering resolution: transcript-level captures isoform-specific states; gene-level provides power for low-abundance transcripts.
    
    \item Neither approach dominates universally. Researchers should consider their biological question when choosing---fine-grained substructure vs statistical robustness.
\end{enumerate}

This robustness analysis confirms that while the strongest signals are cluster-agnostic, subtle isoform switches reveal themselves only under appropriate partitioning schemes. The choice of clustering strategy is not merely technical---it shapes which biology becomes visible.

\subsection*{Threshold Sensitivity: The Tau Sweep}\label{sec:tau-sweep}

The candidate count depends on the contrast threshold $\tau$. To make the choice principled rather than arbitrary, the analysis re-runs the DTU extraction over a fine grid of $\tau$ values (0.005 to 0.3 at step 0.005) and records, at each threshold, the number of accepted events, the unique genes behind them, and how many survive the strictest threshold that still yields at least ten events (the robust core).

Two properties of the sweep matter for interpretation. First, the event count is monotone decreasing in $\tau$---a sanity check that the extraction responds predictably to the threshold. Second, the fraction of accepted events belonging to the robust core (\texttt{frac\_in\_core}) rises toward one as $\tau$ approaches the core threshold; a plateau in this fraction marks the point where further relaxation adds only weak, non-core switches.

Figure~\ref{fig:tau-sweep} maps that search across a landscape of thresholds. The bulk of significant pairs descends like a long, smooth slope---from 2006 at $\tau = 0.005$ to 21 at $\tau = 0.2$---a nearly featureless gradient that reflects the homogeneous cell-line mixture: every isoform pair carries some reciprocal contrast, so the terrain is a dense and continuous spectrum with no natural break. On such a slope any single threshold looks arbitrary, and this is exactly the trap the sweep is built to expose. As the eye climbs the landscape a highland finally appears: at $\tau = 0.25$ ten events, all belonging to the robust core, hold flat through $\tau = 0.26$---a genuine plateau, ground stable enough to stand on rather than a point on the slope. The published threshold $\tau = 0.2$ reproduces the 21 reported candidates and sits just below that highland, so the human result is anchored at the edge of stable ground rather than lost somewhere on the tail. The mouse terrain is different: it descends in steps, 46 events at $\tau = 0.05$ holding at 41 and 36 across wide intervals, because heterogeneous tissue is dominated by a few high-contrast switches with a sparse tail---so its plateaus are broad ledges, not a single summit.

\begin{figure}[htbp]
    \centering
    \includegraphics[width=0.8\textwidth]{arrigoni_tau_sweep.png}
    \caption{Tau sensitivity sweep for the human dataset: DTU events (blue) and robust-core events (orange) versus the minimum DEA contrast threshold $\tau$. A stable plateau appears at $\tau = 0.25$ with a ten-event core; the published $\tau = 0.2$ (dashed) reproduces the 21 reported candidates.}
    \label{fig:tau-sweep}
\end{figure}

The plateau provides an empirical anchor for $\tau$: the reported thresholds (0.2 human, 0.05 mouse) were chosen within or adjacent to stable regions, so the final candidate lists are not an artifact of a single fragile threshold value. The sweep does not, however, replace a permutation-based null distribution; it remains an internal-consistency check on how the accepted set changes with the threshold.

\section{Discussion}\label{sec:discussion}

These findings are now set within their broader methodological and biological context.

\subsection*{Key Findings Interpretation}

Two conclusions follow from these results. The top candidates (Meg3, Malat1, TPM1) align with documented biology---lncRNA switching in neurogenesis and ribosomal protein regulation---which indicates that the algorithm recovers genuine signal rather than technical artifacts. Convergence across a human cell line mixture and mouse brain tissue further indicates that the method holds under different biological contexts rather than a single favourable case study.

\subsection*{Limitations and Caveats}

Despite encouraging results, several constraints temper interpretation:
\begin{itemize}
    \item \textbf{3'-bias invisibility}: Alternative splicing events occurring in transcript mid-regions remain undetectable when only \gls{threePrimeEnd}s are sequenced. The DTU candidates found here represent a subset---those where isoforms differ at their terminal exons or 3' ends.

    \item \textbf{EM rounding effects}: Fractional EM assignments were rounded to integers for COTAN compatibility (Section~\ref{sec:data-filtering}). Low-abundance isoforms with small fractional counts may be distorted by this discretization, potentially introducing false negatives or altered contrast values.

    \item \textbf{Threshold sensitivity}: The choice of $\tau$ lacks formal statistical grounding. The tau sweep (Section~\ref{sec:tau-sweep}) anchors the threshold empirically at a stable plateau, but a permutation-based null distribution would still provide calibrated $p$-values rather than a consistency-based choice.

    \item \textbf{Lack of orthogonal validation}: Without long-read sequencing or targeted RT-PCR\defn{RT-PCR}{Reverse-transcription polymerase chain reaction: a lab technique that experimentally confirms whether a specific isoform is present, used as an orthogonal check on computational DTU calls.} confirmation, false positive rates remain estimated from internal consistency (shared vs exclusive events) rather than ground truth.

This is a limitation common to computational DTU studies but remains critical for biological claims.
\end{itemize}

\subsection*{Future Directions: Extending COTAN for Transcript Analysis}

The framework established here opens several avenues for methodological extension:
\begin{enumerate}
    \item \textbf{Integration with long-read data}: Combining short-read depth (thousands of cells) with PacBio/Nanopore\defn{Long-read sequencing (PacBio, Oxford Nanopore)}{Sequencing technologies that produce full-length transcript reads, giving complete isoform annotation unavailable from short 3'-biased reads.} isoform annotation would anchor DTU calls in complete transcript sequences.

This hybrid approach could validate computational predictions while scaling to population-level sample sizes.

    \item \textbf{Trajectory-aware analysis}: Current clustering assumes discrete cell populations, but differentiation is continuous. Modeling isoform proportions along pseudotime\defn{Pseudotime}{A one-dimensional ordering of cells along a continuous developmental or differentiation trajectory, capturing gradual transitions rather than discrete populations.} trajectories (e.g., using COTAN's framework adapted for regression rather than categorical comparison) could reveal dynamic switching patterns during transitions.

    \item \textbf{Allele-specific expression extension}: COTAN's contingency table approach naturally extends to phased variants. Detecting allele-biased isoform usage would connect regulatory genetics (eQTLs, sQTLs)\defn{eQTL / sQTL}{Genomic variants that measurably affect gene expression (eQTL) or splicing/isoform usage (sQTL); linking them to single-cell isoform switches connects population genetics with regulation.} to splicing outcomes in single cells---bridging population genomics with transcript regulation.

    \item \textbf{Benchmarking against bulk DTU tools}: Systematic comparison with rMATS, SUPPA2, and DRIMSeq \cite{love2018swimming, draper2025selecting} on matched pseudo-bulks would contextualize sensitivity and specificity. This benchmark could identify regimes where single-cell methods add unique value versus when aggregation suffices.
\end{enumerate}

These extensions remain beyond the scope of this thesis, but they indicate that the framework admits further adaptation as new data types and biological questions emerge.
